\documentclass{article}

% ============================================================
% CITATION STYLE
% ============================================================

\PassOptionsToPackage{numbers,sort&compress}{natbib}


% ============================================================
% NEURIPS FORMAT
% ============================================================

\IfFileExists{neurips_2026.sty}{
  \usepackage[preprint]{neurips_2026}
}{
  % Local fallback only. Overleaf should use neurips_2026.sty.
  \usepackage[margin=0.7in]{geometry}
}


% ============================================================
% PACKAGES
% ============================================================

\usepackage[utf8]{inputenc}
\usepackage[T1]{fontenc}

\usepackage{hyperref}
\usepackage{url}

\usepackage{booktabs}
\usepackage{multirow}
\usepackage{array}
\usepackage{tabularx}

\usepackage{amsfonts}
\usepackage{amsmath}
\usepackage{amssymb}

\usepackage{graphicx}
\usepackage{subcaption}
\usepackage{caption}

\usepackage{microtype}

\usepackage{float}
\usepackage{placeins}

\usepackage{xcolor}


% ============================================================
% COMPACT REPORT FORMATTING
% ============================================================

\captionsetup{
    font=footnotesize,
    labelfont=bf,
    skip=3pt
}

\captionsetup[subfigure]{
    font=scriptsize,
    labelfont=normalfont,
    skip=1pt
}

\setlength{\textfloatsep}{5pt plus 1pt minus 1pt}
\setlength{\floatsep}{4pt plus 1pt minus 1pt}
\setlength{\intextsep}{4pt plus 1pt minus 1pt}

\setlength{\abovecaptionskip}{2pt}
\setlength{\belowcaptionskip}{0pt}

\setlength{\parskip}{2pt}

% Compile safely even if a planned figure has not yet been generated.
\newcommand{\figorplaceholder}[2]{%
\IfFileExists{#1}{%
    \includegraphics[width=#2]{#1}%
}{%
    \fbox{\parbox[c][0.18\textheight][c]{0.92\linewidth}{%
        \centering
        Figure not generated yet:\\
        \texttt{\detokenize{#1}}
    }}%
}%
}

% Values that were not serialized by the experiment scripts are reported explicitly.
\newcommand{\notrecorded}{\textit{not saved}}


% ============================================================
% TITLE
% ============================================================

\title{
ATML PA2 Report\\
Preference Optimization, Online RL, Safety, and Feedback Sources
}


% ============================================================
% AUTHOR
% ============================================================

\author{
Muhammad Abdullah Khan\\
Lahore University of Management Sciences
}


% ============================================================
% DOCUMENT
% ============================================================

\begin{document}

\maketitle


% ============================================================
% ABSTRACT
% ============================================================

\begin{abstract}
% WRITE THIS LAST, IN YOUR OWN WORDS.
%
% Suggested content:
% - overall goal of the assignment
% - DPO / PPO / GRPO / safety / RLVR-vs-RLAIF
% - strongest empirical finding
% - important limitation or disagreement between metric and quality
%
% IMPORTANT MANUAL REQUIREMENT:
% End the abstract with the public GitHub repository link.
%
% Repository: \url{YOUR_GITHUB_REPOSITORY_URL}
\end{abstract}

\vspace{3pt}


% ============================================================
% 1. INTRODUCTION
% ============================================================

\section{Introduction}

% WRITE IN YOUR OWN WORDS.
%
% Suggested flow:
% 1. Post-training changes model behavior using preferences/rewards.
% 2. High optimization scores do not necessarily imply better behavior.
% 3. Briefly introduce:
%    - offline DPO
%    - online PPO
%    - critic-free GRPO
%    - safety calibration
%    - exact verifier vs AI feedback
% 4. State the main scientific questions of the assignment.

\vspace{3pt}


% ============================================================
% 2. METHODOLOGY
% ============================================================

\section{Methodology}

\paragraph{Models and training setup.}

% WRITE IN YOUR OWN WORDS.
%
% Mention:
% - Qwen/Qwen2.5-1.5B-Instruct policy family
% - reward model
% - Task 4/5 AI judge
% - LoRA
% - fixed seeds / fixed prompt IDs
% - Apple MPS environment
% - fixed course assets
% - only task-specific ablations changed the designated variable

\begin{center}
\captionof{table}{Core experimental setup and implementation notes.}
\label{tab:setup}
\small
\setlength{\tabcolsep}{5pt}
\begin{tabular}{p{0.24\linewidth}p{0.66\linewidth}}
\toprule
Item & Value / note \\
\midrule
Base policy
& Qwen/Qwen2.5-1.5B-Instruct \\

Reward model
& yavuz-ai/qwen2.5-1.5b-rm-ultrafeedback \\

AI judge
& Qwen/Qwen2.5-3B-Instruct \\

Adapter setup
& LoRA, rank 8, alpha 16, dropout 0.05; Task 1 targets q/v projections \\

Environment
& Apple MPS/fp16 path; frozen reward/judge models use Mac-compatible loading rather than CUDA-specific quantization where necessary \\

Reproducibility
& Fixed course splits, prompt IDs, decoding settings, seeds where practical, and machine-readable saved logs \\

Task 1 preprocessing note
& Prompts exceeding the configured maximum sequence length on their own were skipped rather than truncating the conditioning prompt \\

\bottomrule
\end{tabular}
\end{center}


\paragraph{Evaluation protocol.}

% WRITE IN YOUR OWN WORDS.
%
% Explain:
% - held-out evaluation
% - same prompt/decoding settings within comparisons
% - sampled KL estimator
% - reward-model score is not an absolute quality scale
% - response length uses generated tokens
% - qualitative failure cases complement aggregate metrics
%
% IMPORTANT:
% Include mean + a dispersion statistic for response length
% wherever the manual requests length statistics.


\vspace{3pt}


% ============================================================
% 3. TASK 1
% ============================================================

\section{Task 1: Direct Preference Optimization}

\paragraph{Setup.}

% EXPLAIN IN YOUR OWN WORDS:
%
% - offline chosen/rejected preference pairs
% - trainable policy versus frozen reference
% - LoRA adaptation
% - DPO objective
% - beta
% - standard 1-epoch run
% - short beta forks
% - length-balanced diagnostic run
%
% Task 1 actual config:
% beta = 0.10
% betas = [0.03, 0.10, 0.30]
% learning rate = 2e-5
% batch size = 2
% grad accumulation = 8
% epoch = 1
% max sequence length = 768
% max generation tokens = 256
% LoRA r = 8, alpha = 16, dropout = 0.05
% target modules = q_proj, v_proj


% ------------------------------------------------------------
% TASK 1.1 + 1.2
% ------------------------------------------------------------

\paragraph{Standard DPO and regularization-strength study.}

% INTERPRET:
% - standard run versus the three matched short beta forks
% - standard full run has a DIFFERENT training budget
% - preference accuracy and reward score need not agree
% - relationship with beta was not monotonic in the saved runs

\begin{figure}[H]
\centering
\figorplaceholder{figures/task1_dpo_beta_ablation.png}{0.86\linewidth}
\caption{
Effect of DPO beta on held-out preference accuracy, reward-model score,
and response length. The standard one-epoch model and the short beta forks
use different training budgets and should not be interpreted as a directly
matched comparison.
}
\label{fig:task1-beta}
\end{figure}


\begin{center}
\captionof{table}{Task 1 DPO summary.}
\label{tab:task1-dpo}
\small
\setlength{\tabcolsep}{4pt}
\resizebox{\linewidth}{!}{%
\begin{tabular}{lrrrrrrr}
\toprule
Run & Beta & Steps & Pref. acc. & Loss & KL & Reward & Mean len. (SD) \\
\midrule
standard
& 0.10 & 94 & 0.641 & 0.674 & 0.00037 & 0.896 & 156.4 (106.6) \\

beta\_0\_03
& 0.03 & 38 & 0.631 & 0.687 & -0.00028 & 0.945 & 157.9 (104.9) \\

beta\_0\_10
& 0.10 & 38 & 0.645 & 0.687 & -0.00030 & 0.933 & 153.3 (105.8) \\

beta\_0\_30
& 0.30 & 38 & 0.631 & 0.688 & -0.00023 & 0.954 & 157.5 (105.8) \\
\bottomrule
\end{tabular}}
\end{center}

% REQUIRED LENGTH DISPERSION:
% Add response-length SD or IQR to either:
% - this table,
% - the figure,
% - or one sentence in the text.
%
% Sample SD computed from the original saved response_length fields.


% ------------------------------------------------------------
% TASK 1.3
% ------------------------------------------------------------

\paragraph{Length-confounding study.}

% REQUIRED DISCUSSION:
% - dataset length correlation is NOT the same thing as generated-length shift
% - compare standard vs length-balanced DPO
% - interpret EACH length stratum separately
% - include common word-limit compliance
% - do not claim length explains everything


\begin{figure}[H]
\centering
\figorplaceholder{figures/task1_length_bias_summary.png}{0.82\linewidth}
\caption{
Aggregate comparison of standard and length-balanced DPO on the
length-controlled evaluation.
}
\label{fig:task1-length}
\end{figure}


\begin{center}
\captionof{table}{
Aggregate Task 1 length-balancing comparison on the saved length-stratified evaluation.
}
\label{tab:task1-length}
\small
\setlength{\tabcolsep}{5pt}
\begin{tabular}{lrrrrr}
\toprule
Run & Pref. acc. & Loss & KL & Reward & Mean length (SD) \\
\midrule
standard DPO
& 0.561 & 0.682 & 0.00093 & 1.411 & 142.8 (98.2) \\

length-balanced DPO
& 0.595 & 0.676 & 0.00416 & 1.561 & 140.2 (95.4) \\
\bottomrule
\end{tabular}
\end{center}


% REQUIRED MANUAL EVIDENCE:
% These should be recovered/calculated if possible before submission.

\begin{center}
\captionof{table}{Frozen-checkpoint post-hoc evaluation of the existing adapters:
preference accuracy on 78 preferred-longer, 80 length-matched and 79
rejected-longer pairs; compliance on the fixed ten word-limit prompts.}
\label{tab:task1-strata}
\small
\resizebox{\linewidth}{!}{%
\begin{tabular}{lrrrr}
\toprule
Run & Preferred longer & Length matched & Rejected longer & Word-limit compliance \\
\midrule
standard DPO & 0.4359 & 0.5625 & 0.6835 & 0.7000 \\
length-balanced DPO & 0.4487 & 0.5625 & 0.7722 & 0.6000 \\
\bottomrule
\end{tabular}}
\end{center}


\paragraph{Qualitative evidence.}

% REQUIRED, NOT OPTIONAL.
%
% Include two very short examples:
%
% Example A:
% preference/reward metric improves but actual quality is worse
% due to correctness, concision, or instruction compliance.
%
% Example B:
% illustrates length bias and/or word-limit compliance.
%
% Quote only the minimum response text needed.




% ============================================================
% 4. TASK 2
% ============================================================

\section{Task 2: Proximal Policy Optimization}

\paragraph{Setup.}

% EXPLAIN IN YOUR OWN WORDS:
%
% - continued from supplied PPO midpoint checkpoint
% - online rollout generation
% - reward model
% - learned critic/value model
% - advantage estimation
% - pi_old vs pi_theta
% - clipping epsilon
% - pi_ref and KL pressure
%
% Actual config:
% standard continuation updates = 20
% fork updates = 8
% ppo epochs = 2
% policy LR = 3e-6
% value LoRA LR = 1e-4
% value-head LR = 3e-4
% clip epsilon = 0.20
% clipping ablations = [0.05, 0.20, 0.50]
% KL beta = 0.10
% KL ablations = [0.00, 0.10, 0.20]


% ------------------------------------------------------------
% TASK 2.1
% ------------------------------------------------------------

\paragraph{Standard PPO continuation and training dynamics.}

\begin{figure}[H]
\centering
\figorplaceholder{figures/task2_ppo_standard_trajectory.png}{0.92\linewidth}
\caption{
Standard PPO continuation diagnostics across 20 updates:
reward, sampled KL, policy/value loss, entropy, clip fraction,
gradient norm, and response length.
}
\label{fig:task2-trajectory}
\end{figure}


\begin{center}
\captionof{table}{Standard PPO held-out evaluation and compute summary.}
\label{tab:task2-standard}
\small
\setlength{\tabcolsep}{5pt}
\begin{tabular}{lr}
\toprule
Metric & Value \\
\midrule
Evaluation examples & 200 \\
Reward-model score & 1.4073 \\
Sampled KL & -0.0000587 \\
Entropy & 0.7834 \\
Mean response length & 301.38 \\
Response-length sample SD & 247.22 tokens \\
Termination rate & 0.920 \\
Peak VRAM / peak MPS memory & not recorded; summary reports 0 \\
Wall-clock time & 6.32 min \\
\bottomrule
\end{tabular}
\end{center}


% ------------------------------------------------------------
% TASK 2.2
% ------------------------------------------------------------

\paragraph{Clipping study.}

% EXPLAIN:
% - cached rollout lets epsilon be changed on IDENTICAL data
% - epsilon 0.05 strongly constrains more tokens
% - epsilon 0.20 barely binds
% - epsilon 0.50 effectively does not bind
% - matched short continuation forks test whether this affects behavior

\begin{figure}[H]
\centering
\figorplaceholder{figures/task2_ppo_ablations.png}{0.88\linewidth}
\caption{
PPO ablations: cached clipping diagnostics, matched clipping forks,
and KL-pressure forks.
}
\label{fig:task2-ablations}
\end{figure}


\begin{center}
\captionof{table}{Cached clipping diagnostic.}
\label{tab:task2-clipping-cache}
\small
\setlength{\tabcolsep}{7pt}
\begin{tabular}{rrr}
\toprule
$\epsilon$ & Clip fraction & Affected-token fraction \\
\midrule
0.05 & 0.1373 & 0.1373 \\
0.20 & 0.0018 & 0.0018 \\
0.50 & 0.0000 & 0.0000 \\
\bottomrule
\end{tabular}
\end{center}


\begin{center}
\captionof{table}{Matched PPO clipping and KL-pressure forks.}
\label{tab:task2-ablations}
\small
\setlength{\tabcolsep}{4pt}
\begin{tabular}{lrrrrr}
\toprule
Condition & Reward & KL & Entropy & Mean len. & Max. preclip grad. \\
\midrule
clip 0.05
& 1.260 & 0.000048 & 0.783 & 299.3 & 1.909 \\

clip 0.20
& 1.285 & 0.000038 & 0.786 & 282.1 & 2.568 \\

clip 0.50
& 1.285 & 0.000038 & 0.786 & 282.1 & 2.568 \\

KL 0.00
& 1.292 & -0.000005 & 0.775 & 289.9 & 1.899 \\

KL 0.10
& 1.285 & 0.000038 & 0.786 & 282.1 & 2.568 \\

KL 0.20
& 1.347 & -0.000036 & 0.779 & 308.1 & 2.432 \\
\bottomrule
\end{tabular}
\end{center}

% DEFINE THE STABILITY STATISTIC IN THE TEXT.
% Example possibilities, depending on what is actually saved:
% - standard deviation of reward across fork updates
% - standard deviation of policy loss
% - gradient norm variance
% - maximum gradient norm
%
% Do NOT invent one; compute from saved logs.


% ------------------------------------------------------------
% TASK 2.3
% ------------------------------------------------------------

\paragraph{KL pressure and reward overoptimization.}

% REQUIRED DISCUSSION:
% - beta_KL changes pressure against drifting from pi_ref
% - relationship in your short runs was non-monotonic
% - do not claim beta_KL=0.20 is universally best
% - inspect whether reward gains correspond to actual quality
%
% IMPORTANT CROSS-TASK INSIGHT TO USE LATER IN DISCUSSION:
% KL regularization anchors the policy to the chosen reference,
% but the reference itself may contain undesirable tendencies
% (for example verbosity), so strong KL can preserve these too.


\paragraph{Qualitative evidence.}

% REQUIRED:
% - one case where higher reward corresponds to better response quality
% - one case where reward and response quality disagree




% ============================================================
% 5. TASK 3
% ============================================================

\section{Task 3: Group Relative Policy Optimization}

\paragraph{Setup.}

% EXPLAIN IN YOUR OWN WORDS:
%
% - GRPO removes PPO's critic
% - samples K completions for the same prompt
% - reward-normalized group-relative advantage
% - uninformative groups when reward std ~= 0
% - PPO-style pi_theta/pi_old ratio and clipping remain
% - KL to pi_ref remains
% - standard K = 4
%
% Actual config:
% standard updates = 20
% fork updates = 8
% policy epochs = 1
% LR = 5e-6
% epsilon = 0.20
% KL beta = 0.10
% K = 4
% group sizes = [2, 4, 8]
% max prompt length = 256
% max completion length = 512
% mask truncated completions = true


% ------------------------------------------------------------
% TASK 3.1
% ------------------------------------------------------------

\paragraph{Standard GRPO continuation.}

\begin{figure}[H]
\centering
\figorplaceholder{figures/task3_grpo_standard_trajectory.png}{0.92\linewidth}
\caption{
Standard GRPO continuation diagnostics across training:
reward, sampled KL, within-group reward standard deviation,
uninformative-group rate, policy loss, gradient norm,
entropy, and response length.
}
\label{fig:task3-trajectory}
\end{figure}


\begin{center}
\captionof{table}{Standard GRPO held-out evaluation and compute summary.}
\label{tab:task3-standard}
\small
\setlength{\tabcolsep}{5pt}
\begin{tabular}{lr}
\toprule
Metric & Value \\
\midrule
Evaluation examples & 200 \\
Reward-model score & 1.1125 \\
Sampled KL & 0.0000617 \\
Entropy & 0.8278 \\
Mean response length & 261.12 \\
Response-length sample SD & 187.65 tokens \\
Termination rate & 0.785 \\
Truncation rate & 0.215 \\
Peak VRAM / peak MPS memory & not recorded; summary reports 0 \\
Wall-clock time & 18.37 min \\
\bottomrule
\end{tabular}
\end{center}


% ------------------------------------------------------------
% TASK 3.2
% ------------------------------------------------------------

\paragraph{Group-size study.}

% IMPORTANT:
% Equal TOTAL generation budget.
%
% K=2 -> 96 prompt groups
% K=4 -> 48 prompt groups
% K=8 -> 24 prompt groups
%
% Discuss:
% - informativeness
% - group reward std
% - variance of group-relative signal
% - low/high difficulty bins
% - larger K means more completions per prompt but fewer prompts


\begin{figure}[H]
\centering
\figorplaceholder{figures/task3_group_size_and_normalization.png}{0.80\linewidth}
\caption{
GRPO group-size diagnostics and normalization comparison.
Prompt difficulty bins are defined using the fixed diagnostic rule
used in the analysis script.
}
\label{fig:task3-group-normalization}
\end{figure}


\begin{center}
\captionof{table}{GRPO equal-generation group-size diagnostic.}
\label{tab:task3-groups}
\small
\setlength{\tabcolsep}{3pt}
\begin{tabular}{rrrrrrr}
\toprule
$K$
& Groups
& Info. rate
& Mean group std
& Rel.-signal var.
& Low-bin info.
& High-bin info. \\
\midrule
2
& 96
& 0.9375
& 0.3846
& 0.9375
& 0.8958
& 0.9792 \\

4
& 48
& 0.9583
& 0.5410
& 0.9583
& 0.9167
& 1.0000 \\

8
& 24
& 0.9583
& 0.6191
& 0.9583
& 0.9167
& 1.0000 \\
\bottomrule
\end{tabular}
\end{center}


% ------------------------------------------------------------
% TASK 3.3
% ------------------------------------------------------------

\paragraph{Length-normalization study.}

% EXPLAIN:
% - canonical GRPO normalizes each completion by its own length
% - this changes per-token weighting for short vs long completions
% - Dr.-GRPO changes normalization while keeping reward fixed
% - therefore systematic differences can be attributed to the
%   optimization normalization rather than a new feedback source


\begin{center}
\captionof{table}{Canonical versus Dr.-GRPO matched short continuations.}
\label{tab:task3-normalization}
\small
\setlength{\tabcolsep}{4pt}
\begin{tabular}{lrrrr}
\toprule
Condition
& Reward
& KL
& Mean len.
& Trunc. \\
\midrule
canonical GRPO
& 0.9863
& 0.000005
& 260.5
& 0.225 \\

Dr.-GRPO
& 0.9680
& 0.000055
& 262.9
& 0.270 \\
\bottomrule
\end{tabular}
\end{center}


\begin{center}
\captionof{table}{Post-hoc normalization diagnostic on the fixed cache partitioned into K=4 groups,
with identical rewards and masks. Mean absolute surrogate derivative with
respect to token log probability is evaluated at new=old. Short: at most
206 tokens (62 completions); long: over 206 tokens (60 completions).
Seventy nonterminated or over-512-token completions are masked. These are
log-probability derivatives, not parameter-gradient norms from the K=4 forks.}
\label{tab:task3-length-gradients}
\small
\begin{tabular}{llrr}
\toprule
Normalization & Bin & Mean per-token absolute derivative & Mean completion sum \\
\midrule
GRPO & Short & $1.4834\times10^{-4}$ & 0.0039135 \\
GRPO & Long & $1.1797\times10^{-5}$ & 0.0042074 \\
Dr.-GRPO & Short & $7.6436\times10^{-6}$ & 0.0005264 \\
Dr.-GRPO & Long & $8.2176\times10^{-6}$ & 0.0030643 \\
\bottomrule
\end{tabular}
\end{center}

\paragraph{Qualitative evidence.}

% REQUIRED:
% Include a brief example useful for interpreting either:
% - an informative/uninformative group
% - or the length-normalization behavior




% ============================================================
% 6. TASK 4
% ============================================================

\section{Task 4: Safety Calibration}

\paragraph{Setup.}

% EXPLAIN IN YOUR OWN WORDS:
%
% - fixed XSTest prompts
% - SFT / standard DPO / standard PPO / standard GRPO
% - deterministic generation
% - fixed categorical AI judge
% - safe vs unsafe prompt distinction
% - safety is calibration, not simply "refuse more"


% ------------------------------------------------------------
% TASK 4.1 + 4.2
% ------------------------------------------------------------

\paragraph{Policy-level safety calibration.}

\begin{figure}[H]
\centering
\figorplaceholder{figures/task4_safety_policy_summary.png}{0.78\linewidth}
\caption{
Policy-level XSTest safety calibration across SFT, DPO, PPO, and GRPO.
}
\label{fig:task4-summary}
\end{figure}


\begin{center}
\captionof{table}{Task 4 safety calibration summary from the fixed AI judge.}
\label{tab:task4-safety}
\small
\setlength{\tabcolsep}{3pt}
\begin{tabular}{lrrrrrrr}
\toprule
Policy
& Safe ans.
& Over-ref.
& Just. ref.
& Unsafe comp.
& Ambig.
& Calibrated
& Mean tok. \\
\midrule
SFT
& 0.704 & 0.272 & 0.805 & 0.190 & 0.011 & 0.749 & 105.9 \\

DPO
& 0.672 & 0.316 & 0.805 & 0.195 & 0.002 & 0.731 & 109.4 \\

PPO
& 0.704 & 0.272 & 0.805 & 0.190 & 0.011 & 0.749 & 106.6 \\

GRPO
& 0.708 & 0.264 & 0.810 & 0.185 & 0.013 & 0.753 & 107.1 \\
\bottomrule
\end{tabular}
\end{center}


\paragraph{Category-level behavior.}

% MANUAL REQUIRES THE JUDGE-LABEL DISTRIBUTION ACROSS CATEGORIES
% FOR ALL FOUR POLICIES.
%
% The existing over-refusal / unsafe-compliance heatmaps are useful,
% but a fuller distribution figure/table should also exist.

\begin{figure}[H]
\centering
\figorplaceholder{figures/task4_category_label_distribution.png}{0.84\linewidth}
\caption{
Distribution of AI-judge behavior labels across XSTest categories
for SFT, DPO, PPO, and GRPO.
}
\label{fig:task4-category-distribution}
\end{figure}


% Existing error-focused heatmap can be retained if space allows.

\begin{figure}[H]
\centering
\figorplaceholder{figures/task4_category_heatmaps.png}{0.78\linewidth}
\caption{
Category-level error rates. Left: over-refusal on safe prompts.
Right: unsafe compliance on unsafe prompts.
}
\label{fig:task4-category-errors}
\end{figure}


% ------------------------------------------------------------
% TASK 4.3
% ------------------------------------------------------------

\paragraph{Manual audit.}

% REQUIRED:
% - manually label the fixed 60-example subset WITHOUT viewing AI labels first
% - same five behavior labels
% - automated-vs-manual agreement
% - confusion-style breakdown
% - ambiguous information
%
% manual_audit_ids.csv has 60 example IDs
% manual_audit_sheet.csv has 240 policy-response rows


\begin{center}
\captionof{table}{AI-judge versus manual-audit agreement.}
\label{tab:task4-manual}
\small
\setlength{\tabcolsep}{5pt}
\begin{tabular}{lr}
\toprule
Metric & Value \\
\midrule
Overall agreement & 0.6333 \\
AI ambiguous rate & 0.0000 \\
Manual ambiguous rate & 0.0000 \\
Audited prompt IDs & 60 \\
Audited policy-response pairs & 240 \\
\bottomrule
\end{tabular}
\end{center}


\begin{figure}[H]
\centering
\figorplaceholder{figures/task4_manual_audit_confusion.png}{0.58\linewidth}
\caption{
Confusion-style comparison between the fixed AI judge and manual labels
on the required audit subset.
}
\label{fig:task4-manual-confusion}
\end{figure}


\paragraph{Qualitative disagreement cases.}

% REQUIRED:
% Include short examples covering:
%
% 1. harmful compliance
% 2. justified refusal
% 3. exaggerated refusal / over-refusal
%
% For each, state whether the issue reflects:
% - policy behavior,
% - judge error,
% - or both.




% ============================================================
% 7. TASK 5
% ============================================================

\section{Task 5: Feedback Source --- RLVR versus RLAIF}

\paragraph{Setup.}

% EXPLAIN IN YOUR OWN WORDS:
%
% - SFT / supplied RLVR / supplied RLAIF policies
% - same policy family
% - group-based optimization broadly matched
% - RLVR uses exact final-answer verification
% - RLAIF uses frozen AI pairwise preferences
% - GSM8K is in-domain
% - SVAMP is transfer/out-of-domain
% - controlled diagnostic set isolates reward-source behavior


% ------------------------------------------------------------
% TASK 5.1 + 5.3
% ------------------------------------------------------------

\paragraph{In-domain and transfer performance.}

\begin{figure}[H]
\centering
\figorplaceholder{figures/task5_gsm_svamp_policy_comparison.png}{0.82\linewidth}
\caption{
In-domain GSM8K and out-of-domain SVAMP performance for SFT, RLVR,
and RLAIF.
}
\label{fig:task5-policy-comparison}
\end{figure}


\begin{center}
\captionof{table}{
Task 5 exact-answer accuracy, format compliance, and response length.
}
\label{tab:task5-metrics}
\small
\setlength{\tabcolsep}{4pt}
\begin{tabular}{llrrrr}
\toprule
Dataset & Policy & Exact acc. & Format & Mean tok. (SD) & Trunc. \\
\midrule
GSM8K & SFT
& 0.293 & 0.400 & 284.1 (100.4) & 0.047 \\

GSM8K & RLVR
& 0.290 & 0.420 & 283.9 (99.1) & 0.043 \\

GSM8K & RLAIF
& 0.297 & 0.407 & 289.1 (99.3) & 0.043 \\

SVAMP & SFT
& 0.560 & 0.670 & 162.5 (71.5) & 0.010 \\

SVAMP & RLVR
& 0.450 & 0.570 & 168.1 (102.2) & 0.010 \\

SVAMP & RLAIF
& 0.570 & 0.630 & 171.2 (76.0) & 0.000 \\
\bottomrule
\end{tabular}
\end{center}


% REQUIRED:
% Add response-length SD or IQR for these conditions if available.


\begin{center}
\captionof{table}{
Verifier--judge agreement and its change from GSM8K to SVAMP.
}
\label{tab:task5-agreement-transfer}
\small
\setlength{\tabcolsep}{4pt}
\begin{tabular}{lrrr}
\toprule
Pairwise comparison
& GSM8K verifier--judge agreement
& SVAMP verifier--judge agreement
& SVAMP minus GSM8K \\
\midrule
RLVR vs RLAIF
& 0.373 & 0.260 & -0.113 \\

SFT vs RLAIF
& 0.383 & 0.300 & -0.083 \\

SFT vs RLVR
& 0.387 & 0.300 & -0.087 \\
\bottomrule
\end{tabular}
\end{center}

% NOTE:
% Depending on how agreement is defined in the saved analysis,
% it may be more appropriate to report agreement per PAIRWISE COMPARISON
% rather than per policy. Match the implementation's actual definition.
%
% The manual explicitly requires verifier--judge agreement and
% the transfer drop/change where applicable.


% ------------------------------------------------------------
% TASK 5 PAIRWISE RESULTS
% ------------------------------------------------------------

\begin{center}
\captionof{table}{Task 5 pairwise AI-judge outcomes.}
\label{tab:task5-pairwise}
\small
\setlength{\tabcolsep}{4pt}
\begin{tabular}{llrrr}
\toprule
Dataset & Comparison & First wins & Second wins & Ties \\
\midrule
GSM8K
& RLVR vs RLAIF
& 103 & 95 & 102 \\

GSM8K
& SFT vs RLAIF
& 99 & 103 & 98 \\

GSM8K
& SFT vs RLVR
& 98 & 105 & 97 \\

SVAMP
& RLVR vs RLAIF
& 37 & 37 & 26 \\

SVAMP
& SFT vs RLAIF
& 44 & 36 & 20 \\

SVAMP
& SFT vs RLVR
& 34 & 40 & 26 \\
\bottomrule
\end{tabular}
\end{center}


% ------------------------------------------------------------
% TASK 5.2
% ------------------------------------------------------------

\paragraph{Controlled feedback diagnostics.}

% REQUIRED:
% For EACH perturbation category and EACH feedback mechanism:
% - intended-better rate
% - tie rate
% - wrong-preference rate
%
% Also report:
% - S_reason
% - S_outcome

\begin{figure}[H]
\centering
\figorplaceholder{figures/task5_controlled_feedback_diagnostics.png}{0.78\linewidth}
\caption{
Controlled reward-source diagnostic. Bars show how often each feedback
mechanism prefers the diagnostically intended-better response, ties,
or prefers the intended-worse response.
}
\label{fig:task5-controlled}
\end{figure}


\begin{center}
\captionof{table}{
Controlled diagnostic preference rates for the exact verifier and AI judge.
}
\label{tab:task5-controlled}
\small
\setlength{\tabcolsep}{3pt}
\begin{tabular}{llrrr}
\toprule
Feedback
& Perturbation
& Better
& Tie
& Wrong \\
\midrule

Exact verifier
& distractor robustness
& 1.00 & 0.00 & 0.00 \\

Exact verifier
& filler susceptibility
& 0.00 & 1.00 & 0.00 \\

Exact verifier
& final vs reasoning
& 1.00 & 0.00 & 0.00 \\

Exact verifier
& outcome sensitivity
& 1.00 & 0.00 & 0.00 \\

Exact verifier
& reasoning sensitivity
& 0.00 & 1.00 & 0.00 \\

\midrule

AI judge
& distractor robustness
& 0.00 & 1.00 & 0.00 \\

AI judge
& filler susceptibility
& 0.10 & 0.35 & 0.55 \\

AI judge
& final vs reasoning
& 0.20 & 0.80 & 0.00 \\

AI judge
& outcome sensitivity
& 0.15 & 0.85 & 0.00 \\

AI judge
& reasoning sensitivity
& 0.10 & 0.60 & 0.30 \\

\bottomrule
\end{tabular}
\end{center}


\begin{center}
\captionof{table}{
Controlled reasoning and outcome sensitivity summary.
}
\label{tab:task5-sensitivity}
\small
\setlength{\tabcolsep}{7pt}
\begin{tabular}{lrr}
\toprule
Feedback mechanism
& $S_{\mathrm{reason}}$
& $S_{\mathrm{outcome}}$ \\
\midrule
Exact verifier
& 0.00 & 1.00 \\

AI judge
& 0.10 & 0.15 \\
\bottomrule
\end{tabular}
\end{center}


\paragraph{Qualitative diagnostic cases.}

% REQUIRED:
% Include short examples spanning:
%
% 1. reasoning corruption while final answer remains correct
% 2. persuasive / irrelevant filler
% 3. wrong final answer with gold-number distractor
%
% Use them to explain verifier--judge disagreement.


\paragraph{Feedback-source comparison.}

% REQUIRED DISCUSSION:
%
% Compare RLVR and RLAIF in terms of:
% - feedback coverage
% - noise
% - exploitability
% - inference cost
%
% Do NOT compare raw binary RLVR reward numerically against
% RLAIF pairwise win-rate as though they share a calibrated scale.




% ============================================================
% 8. CROSS-TASK DISCUSSION
% ============================================================

\section{Cross-Task Discussion}

% WRITE IN YOUR OWN WORDS.
%
% Suggested structure:
%
% 1. Preference strength versus policy drift
%    - DPO beta
%    - PPO KL pressure
%    - GRPO normalization
%    - KL / reward / entropy / length
%
% 2. Offline versus online feedback
%    - DPO fixed preference pairs
%    - PPO/GRPO sampled rollouts
%
% 3. Optimization bias
%    - DPO length confounding
%    - PPO clipping / reward overoptimization
%    - GRPO group informativeness / sequence normalization
%
% 4. Safety calibration
%    - whether better preference/reward metrics implied better safety
%    - safe over-refusal versus unsafe compliance
%
% 5. Reward-source design
%    - learned preference reward
%    - exact verification
%    - AI feedback
%
% IMPORTANT KL DISCUSSION POINT:
%
% KL regularization protects useful behavior by anchoring the optimized
% policy to the reference policy. However, the reference policy is not
% necessarily ideal or unbiased. If it already contains undesirable
% tendencies such as verbosity, excessive caution, or other stylistic
% biases, stronger KL pressure can also preserve those tendencies.
%
% This means "lower KL" should not automatically be interpreted as
% "better alignment."




% ============================================================
% 9. LIMITATIONS
% ============================================================

\section{Limitations}

% WRITE IN YOUR OWN WORDS.
%
% Potential items supported by this project:
%
% - Apple MPS/fp16 environment rather than the intended CUDA path
% - Mac-compatible unquantized loading for some frozen models
% - Task 1 overlong-prompt filtering
% - small / short ablation budgets
% - sampled KL estimator can be slightly negative
% - learned reward model is not calibrated model quality
% - PPO critic is imperfect
% - GRPO truncation is nontrivial
% - AI safety judge is imperfect
% - strict Task 5 final-answer formatting affects exact accuracy
% - manual audit sample is finite
%
% Do not describe required evidence as "missing" in the final report
% if it can still be recovered before submission.




% ============================================================
% 10. CONCLUSION
% ============================================================

\section{Conclusion}

% WRITE 3--5 SENTENCES IN YOUR OWN WORDS.
%
% Suggested content:
% - strongest empirical conclusion
% - one result showing metrics/objectives can disagree
% - one limitation
% - general lesson about post-training objectives and feedback


% ============================================================
% REFERENCES
% ============================================================

% References do not count toward the 8-page main-content limit.
%
% Recommended citations from the assignment:
% - Rafailov et al. (2023) DPO
% - Schulman et al. (2017) PPO
% - Ouyang et al. (2022) RLHF / InstructGPT
% - Shao et al. (2024) DeepSeekMath / GRPO
% - Liu et al. (2025) GRPO / Dr.-GRPO analysis
% - Rottger et al. (2023) XSTest
% - DeepSeek-R1 paper for RLVR context
% - Lee et al. (2024) RLAIF vs RLHF
%
% If references.bib exists, use it. Otherwise the document still compiles.

\IfFileExists{references.bib}{
    \bibliographystyle{plainnat}
    \bibliography{references}
}{
    \section*{References}
    % Create references.bib before final submission.
}


\end{document}
